% GENERATED FILE. Edit canonical lessons, then run npm run generate:books.
% canonical-source-hash: fnv1a64:15d088f80bf0d2f3
% canonical-lessons: TA-C382-nipantanai, TA-C382-utpatu, TA-C382-takuti, TA-C382-vitivilakku, TA-C382-ventumanal

\chapter{On One Condition}
\label{ch:ta-on-one-condition}
\hlchaptermodality{382}

\begin{chapteropening}
\textbf{By the end of this chapter:} \emph{I can set a condition and say what is subject to it.}

Complete the last lesson of chapter 382: I can set a condition and say what is subject to it.
\end{chapteropening}

\section[\texorpdfstring{\ta{நிபந்தனை} (nipantaṉai)}{நிபந்தனை (nipantaṉai)}]{\ta{நிபந்தனை} (nipantaṉai) — a condition}

\label{lesson:TA-C382-nipantanai}



\begin{quote}
Before the new one: say the Tamil for to unite, then the Tamil for a representative.
\end{quote}

\subsection*{You'll want to know: \ta{நிபந்தனை}}
\textbf{\ta{நிபந்தனை}} — \emph{nipantaṉai} — \textquotedblleft{}a condition\textquotedblright{}.

\begin{quote}
\ta{ஒரு} \ta{நிபந்தனையின்} \ta{பேரில்} \ta{நான்} \ta{வருகிறேன்}.
\end{quote}

\begin{quote}
\emph{oru nipantaṉaiyiṉ pēril nāṉ varukiṟēṉ.}
\end{quote}

\begin{quote}
I will come on one condition.
\end{quote}

Plural: \ta{நிபந்தனைகள்}, \emph{nipantaṉaikaḷ}.

\ta{நிபந்தனையின்} \ta{பேரில்}, on condition; \ta{நிபந்தனையின்றி}, unconditionally.

\begin{cousinweb}
\textbf{\ta{நிபந்தனை}} (\emph{nipantaṉai}) is borrowed from Sanskrit \emph{nibandhana}, \textquotedblleft{}tying, binding; a band, a fetter\textquotedblright{} (Monier-Williams): what an agreement is tied to.
\end{cousinweb}

\begin{grammarlens}[title={What you've built}]
One more word for saying what happens if.
\end{grammarlens}

\subsection*{Your turn}
\begin{itemize}
\raggedright
  \item \emph{Say it:} \emph{nipantaṉai}
  \item \emph{Say it:} \emph{nipantaṉai}, once more
  \item \emph{Say it:} \emph{oru nipantaṉaiyiṉ pēril nāṉ varukiṟēṉ.} — I will come on one condition.
  \item \emph{Recall:} say the Tamil for to unite, then the Tamil for a representative, then say \emph{nipantaṉai} again
\end{itemize}

\begin{tcolorbox}[breakable,colback=teal!4,colframe=teal!35!black,title={Before you move on}]
What does \ta{நிபந்தனை} mean? (\textquotedblleft{}A condition\textquotedblright{}.) Say it once more.
\end{tcolorbox}
\section[\texorpdfstring{\ta{உட்படு} (uṭpaṭu)}{உட்படு (uṭpaṭu)}]{\ta{உட்படு} (uṭpaṭu) — to be subject to; to come under}

\label{lesson:TA-C382-utpatu}



\begin{quote}
Before the new one: say the Tamil for a representative, then the Tamil for a condition.
\end{quote}

\subsection*{You'll want to know: \ta{உட்படு}}
\textbf{\ta{உட்படு}} — \emph{uṭpaṭu} — \textquotedblleft{}to be subject to; to come under\textquotedblright{}.

\begin{quote}
\ta{இந்த} \ta{அனுமதி} \ta{சில} \ta{நிபந்தனைகளுக்கு} \ta{உட்பட்டது}.
\end{quote}

\begin{quote}
\emph{inta aṉumati cila nipantaṉaikaḷukku uṭpaṭṭatu.}
\end{quote}

\begin{quote}
This permission is subject to some conditions.
\end{quote}

Past (I …): \ta{உட்பட்டேன்}, \emph{uṭpaṭṭēṉ}.

With -\ta{க்கு}: \ta{விதிகளுக்கு} \ta{உட்பட்டு}, subject to the rules. \ta{உட்பட}, including.

\begin{cousinweb}
\textbf{\ta{உட்படு}} (\emph{uṭpaṭu}) is native: \emph{uḷ}, \textquotedblleft{}inside\textquotedblright{} (DEDR 698), whose ḷ becomes \ta{ட்} before p, and \emph{paṭu}, \textquotedblleft{}to undergo\textquotedblright{}: to come inside a rule.
\end{cousinweb}

\begin{grammarlens}[title={What you've built}]
One more word for saying what happens if.
\end{grammarlens}

\subsection*{Your turn}
\begin{itemize}
\raggedright
  \item \emph{Say it:} \emph{uṭpaṭu}
  \item \emph{Say it:} \emph{uṭpaṭu}, once more
  \item \emph{Say it:} \emph{inta aṉumati cila nipantaṉaikaḷukku uṭpaṭṭatu.} — This permission is subject to some conditions.
  \item \emph{Recall:} say the Tamil for a representative, then the Tamil for a condition, then say \emph{uṭpaṭu} again
\end{itemize}

\begin{tcolorbox}[breakable,colback=teal!4,colframe=teal!35!black,title={Before you move on}]
What does \ta{உட்படு} mean? (\textquotedblleft{}To be subject to; to come under\textquotedblright{}.) Say it once more.
\end{tcolorbox}
\section[\texorpdfstring{\ta{தகுதி} (takuti)}{தகுதி (takuti)}]{\ta{தகுதி} (takuti) — eligibility; a qualification}

\label{lesson:TA-C382-takuti}



\begin{quote}
Before the new one: say the Tamil for a condition, then the Tamil for to be subject to.
\end{quote}

\subsection*{You'll want to know: \ta{தகுதி}}
\textbf{\ta{தகுதி}} — \emph{takuti} — \textquotedblleft{}eligibility; a qualification\textquotedblright{}.

\begin{quote}
\ta{இந்த} \ta{வேலைக்கு} \ta{உங்களுக்குத்} \ta{தகுதி} \ta{இருந்தால்} \ta{விண்ணப்பிக்கலாம்}.
\end{quote}

\begin{quote}
\emph{inta vēlaikku uṅkaḷukkut takuti iruntāl viṇṇappikkalām.}
\end{quote}

\begin{quote}
If you are qualified for this job, you may apply.
\end{quote}

X-\ta{க்குத்} \ta{தகுதி}, eligibility for X. \ta{தகுதியான}, deserving, suitable.

\begin{cousinweb}
\textbf{\ta{தகுதி}} (\emph{takuti}) is native Dravidian: DEDR 3005 lists Tamil \emph{taku}, \textquotedblleft{}to be fit, proper, worthy, adequate\textquotedblright{}, and \emph{takuti}, \textquotedblleft{}fitness, worth\textquotedblright{}.
\end{cousinweb}

\begin{grammarlens}[title={What you've built}]
One more word for saying what happens if.
\end{grammarlens}

\subsection*{Your turn}
\begin{itemize}
\raggedright
  \item \emph{Say it:} \emph{takuti}
  \item \emph{Say it:} \emph{takuti}, once more
  \item \emph{Say it:} \emph{inta vēlaikku uṅkaḷukkut takuti iruntāl viṇṇappikkalām.} — If you are qualified for this job, you may apply.
  \item \emph{Recall:} say the Tamil for a condition, then the Tamil for to be subject to, then say \emph{takuti} again
\end{itemize}

\begin{tcolorbox}[breakable,colback=teal!4,colframe=teal!35!black,title={Before you move on}]
What does \ta{தகுதி} mean? (\textquotedblleft{}Eligibility; a qualification\textquotedblright{}.) Say it once more.
\end{tcolorbox}
\section[\texorpdfstring{\ta{விதிவிலக்கு} (vitivilakku)}{விதிவிலக்கு (vitivilakku)}]{\ta{விதிவிலக்கு} (vitivilakku) — an exception}

\label{lesson:TA-C382-vitivilakku}



\begin{quote}
Before the new one: say the Tamil for to be subject to, then the Tamil for eligibility.
\end{quote}

\subsection*{You'll want to know: \ta{விதிவிலக்கு}}
\textbf{\ta{விதிவிலக்கு}} — \emph{vitivilakku} — \textquotedblleft{}an exception\textquotedblright{}.

\begin{quote}
\ta{இந்த} \ta{விதிக்கு} \ta{ஒரு} \ta{விதிவிலக்கு} \ta{இருக்கிறது}.
\end{quote}

\begin{quote}
\emph{inta vitikku oru vitivilakku irukkiṟatu.}
\end{quote}

\begin{quote}
There is one exception to this rule.
\end{quote}

\ta{விதி}, rule + \ta{விலக்கு}, a setting aside.

\begin{cousinweb}
\textbf{\ta{விதிவிலக்கு}} (\emph{vitivilakku}) is built from \emph{viti}, \textquotedblleft{}rule\textquotedblright{}, borrowed from Sanskrit \emph{vidhi} (CDIAL 11755), and native \emph{vilakku}, \textquotedblleft{}to turn aside, avert\textquotedblright{} (DEDR 5423): a rule turned aside.
\end{cousinweb}

\begin{grammarlens}[title={What you've built}]
One more word for saying what happens if.
\end{grammarlens}

\subsection*{Your turn}
\begin{itemize}
\raggedright
  \item \emph{Say it:} \emph{vitivilakku}
  \item \emph{Say it:} \emph{vitivilakku}, once more
  \item \emph{Say it:} \emph{inta vitikku oru vitivilakku irukkiṟatu.} — There is one exception to this rule.
  \item \emph{Recall:} say the Tamil for to be subject to, then the Tamil for eligibility, then say \emph{vitivilakku} again
\end{itemize}

\begin{tcolorbox}[breakable,colback=teal!4,colframe=teal!35!black,title={Before you move on}]
What does \ta{விதிவிலக்கு} mean? (\textquotedblleft{}An exception\textquotedblright{}.) Say it once more.
\end{tcolorbox}
\section[\texorpdfstring{\ta{வேண்டுமானால்} (vēṇṭumāṉāl)}{வேண்டுமானால் (vēṇṭumāṉāl)}]{\ta{வேண்டுமானால்} (vēṇṭumāṉāl) — if need be; if you like}

\label{lesson:TA-C382-ventumanal}



\begin{quote}
Before the new one: say the Tamil for eligibility, then the Tamil for an exception.
\end{quote}

\subsection*{You'll want to know: \ta{வேண்டுமானால்}}
\textbf{\ta{வேண்டுமானால்}} — \emph{vēṇṭumāṉāl} — \textquotedblleft{}if need be; if you like\textquotedblright{}.

\begin{quote}
\ta{வேண்டுமானால்} \ta{நீ} \ta{என்} \ta{கைப்பேசியைப்} \ta{பயன்படுத்தலாம்}.
\end{quote}

\begin{quote}
\emph{vēṇṭumāṉāl nī eṉ kaippēciyaip payaṉpaṭuttalām.}
\end{quote}

\begin{quote}
If you like, you can use my phone.
\end{quote}

\ta{வேண்டும்}, want, + \ta{ஆனால்}, if it is: the conditional of \ta{வேண்டும்}.

\begin{cousinweb}
\textbf{\ta{வேண்டுமானால்}} (\emph{vēṇṭumāṉāl}) is native: \emph{vēṇṭum}, \textquotedblleft{}it is wanted\textquotedblright{}, from \emph{vēṇṭu} (DEDR 5528), and \emph{āṉāl}, \textquotedblleft{}if it becomes\textquotedblright{}, the conditional of \emph{āku} (DEDR 333).
\end{cousinweb}

\begin{grammarlens}[title={What you've built}]
One more word for saying what happens if.
\end{grammarlens}

\subsection*{Your turn}
\begin{itemize}
\raggedright
  \item \emph{Say it:} \emph{vēṇṭumāṉāl}
  \item \emph{Say it:} \emph{vēṇṭumāṉāl}, once more
  \item \emph{Say it:} \emph{vēṇṭumāṉāl nī eṉ kaippēciyaip payaṉpaṭuttalām.} — If you like, you can use my phone.
  \item \emph{Recall:} say the Tamil for eligibility, then the Tamil for an exception, then say \emph{vēṇṭumāṉāl} again
\end{itemize}

\begin{tcolorbox}[breakable,colback=teal!4,colframe=teal!35!black,title={Before you move on}]
What does \ta{வேண்டுமானால்} mean? (\textquotedblleft{}If need be; if you like\textquotedblright{}.) Say it once more.
\end{tcolorbox}
